\documentclass[11pt]{article}
\usepackage[margin=1in]{geometry}
\usepackage{amsmath,amssymb}
\usepackage[T1]{fontenc}
\usepackage[hidelinks]{hyperref}

\newcommand{\BPBp}{\mathrm{BPBp}}
\newcommand{\R}{\mathbb{R}}

\title{Literature-Implied Status Note for arXiv:1806.09366}
\author{Run fa\_banach\_001, agent\_lane\_01}
\date{2026-06-14}

\begin{document}
\maketitle

\section*{Status}

\textbf{literature\_implied\_answer (partial subcase).}

This note records a known-theorem implication for one scalar-coordinate
reading of the reciprocal Bishop--Phelps--Bollobas question in
Choi--Dantas--Jung--Martin, \emph{The Bishop-Phelps-Bollobas property
and absolute sums}, arXiv:1806.09366.

\section*{Source Question}

In Section 3, after proving inheritance results for absolute summands in the
domain, the source paper discusses reciprocal results.  It observes that
although each pair $(\R,Y)$ has the BPBp trivially, there are Banach spaces
$Y$ such that $(\R\oplus_1 \R,Y)$ does not have the BPBp; since
$\R\oplus_1\R$ is linearly isometric to $\R\oplus_\infty\R$, the same example
also treats the $\ell_\infty$ scalar sum.  The source then asks what happens
for $\ell_p$-sums with $1<p<\infty$.

\section*{Supporting Literature}

Kim--Lee, \emph{Uniform convexity and the Bishop-Phelps-Bollobas property},
Canad. J. Math. 66 (2014), Theorem 3.1, proves that if the domain space
$X$ is uniformly convex, then $(X,Y)$ has the BPBp for every Banach space
$Y$.  The same theorem is quoted in Dantas--Kadets--Kim--Lee--Martin,
arXiv:1709.00032, Introduction, in the form: among known BPBp results,
$(X,Y)$ has the BPBp when $X$ is uniformly convex, citing
Kim--Lee Theorem 3.1 and Acosta--Becerra Guerrero--Mero\~no Theorem 2.2.

\section*{Identification}

For every $1<p<\infty$, the two-dimensional space
\[
  \R\oplus_p\R = \ell_p^2
\]
is uniformly convex.  Therefore Kim--Lee's uniformly convex-domain theorem
applies directly and gives that
\[
  (\ell_p^2,Y)
\]
has the BPBp for every Banach space $Y$.

Consequently, the specific scalar-coordinate obstruction used in the source
for $p=1$ and $p=\infty$ cannot occur for $1<p<\infty$: even though the
coordinate pairs $(\R,Y)$ are only trivially BPBp, the full scalar
$\ell_p^2$ domain is automatically BPBp with every range.

\section*{Scope}

This is not a new proof and not a full solution to the arbitrary reciprocal
problem for $X_1\oplus_p X_2$.  It answers only the scalar two-summand subcase
$X_1=X_2=\R$ that is naturally suggested by the source paper's displayed
$\R\oplus_1\R$ and $\R\oplus_\infty\R$ counterexample discussion.  The
general implication
\[
  (X_1,Y), (X_2,Y) \text{ have BPBp } \Longrightarrow
  (X_1\oplus_p X_2,Y) \text{ has BPBp}
\]
for arbitrary Banach spaces remains outside the scope of this packet.

\section*{Provenance}

Local evidence used:
\begin{itemize}
  \item arXiv:1806.09366 source, Section 3, line 451 in the parsed TeX.
  \item arXiv:1709.00032 source, Introduction, quoting the uniformly
  convex-domain theorem and the $\ell_1^2$ obstruction context.
  \item arXiv:1605.00245 source, Example after Proposition 3.??, quoting
  Kim--Lee Corollary 3.3 for the two-dimensional real characterization.
\end{itemize}

\begin{thebibliography}{9}

\bibitem{CDJM}
Y. S. Choi, S. Dantas, M. Jung, and M. Martin,
\emph{The Bishop-Phelps-Bollobas property and absolute sums},
arXiv:1806.09366.

\bibitem{KL}
S. K. Kim and H. J. Lee,
\emph{Uniform convexity and the Bishop-Phelps-Bollobas property},
Canad. J. Math. \textbf{66} (2014), 373--386.

\bibitem{DKKLM}
S. Dantas, V. Kadets, S. K. Kim, H. J. Lee, and M. Martin,
\emph{There is no operatorwise version of the Bishop-Phelps-Bollobas property},
arXiv:1709.00032.

\end{thebibliography}

\end{document}
